\documentclass[11pt,a4paper]{article}
\usepackage[utf8]{inputenc}
\usepackage[vietnamese]{babel}
\usepackage{amsmath,amssymb}
\usepackage{geometry}
\geometry{a4paper, total={180mm,260mm}, left=15mm, right=15mm, top=18mm, bottom=18mm}
\usepackage{tikz}
\usetikzlibrary{shapes.geometric, arrows.meta, positioning, calc}
\usepackage{tcolorbox}
\tcbuselibrary{skins,breakable}
\usepackage{xcolor}
\usepackage{multicol}
\usepackage{array}
\usepackage{fontawesome}
\usepackage{fancyhdr}

\pagestyle{fancy}
\fancyhf{}
\renewcommand{\headrulewidth}{0.5pt}
\renewcommand{\footrulewidth}{0.5pt}
\fancyhead[L]{\small\bfseries\color{teal}\faBook\ SỔ TAY KIẾN THỨC TOÁN 3}
\fancyhead[R]{\small\bfseries\color{teal}Lớp Toán Thầy Lĩnh -- loptoanthaylinh.vn}
\fancyfoot[C]{\small\color{gray}-- Sổ tay kiến thức: Bài 75. Thực hành trải nghiệm thống kê --}

\definecolor{primary}{HTML}{005F73}
\definecolor{secondary}{HTML}{0A9396}
\definecolor{accent}{HTML}{EE9B00}
\definecolor{danger}{HTML}{AE2012}
\definecolor{boxbg}{HTML}{F4F6F8}

\begin{document}

% --- BANNER TIÊU ĐỀ ---
\begin{tcolorbox}[
  enhanced,
  colback=white,
  colframe=primary,
  arc=3mm,
  boxrule=1.2pt,
  left=4mm, right=4mm, top=2.5mm, bottom=2.5mm,
  drop shadow=black!10!white
]
  \centering
  {\Large\bfseries\color{primary}\faBookmark\ BÀI 75: THỰC HÀNH TRẢI NGHIỆM THỐNG KÊ}\\[1.5mm]
  {\small\itshape\color{black!75}Kĩ năng lập bảng kiểm đếm thực tế, thống kê số lần xuất hiện và rút ra nhận xét xác suất trên 1 trang A4}
\end{tcolorbox}

\vspace{-2mm}


\noindent
\begin{minipage}[t]{0.49\textwidth}\vspace{0pt}
\begin{tcolorbox}[
  enhanced,
  colback=blue!4!white,
  colframe=blue!75!black,
  arc=2mm, boxrule=0.8pt,
  left=2mm, right=2mm, top=1.8mm, bottom=1.8mm,
  title={\footnotesize\bfseries\color{white}\faKey\ 1. Kĩ thuật kiểm đếm bằng vạch}
]
{\footnotesize
\begin{itemize}\setlength{\itemsep}{1.5pt}\setlength{\leftskip}{-2.5mm}
  \item Khi làm thí nghiệm (gieo xúc xắc, quay vòng quay), mỗi lần xuất hiện một kết quả ta ghi $1$ gạch ($|$).
  \item Cứ $5$ gạch ta nhóm lại thành một cụm để dễ đếm: bốn gạch đứng và một gạch chéo.
  \item \textbf{Tổng số lần thử nghiệm:}
  $$\text{Tổng số lần} = \text{Tổng số vạch kiểm đếm}$$
  \item Lập bảng tổng kết số lần xuất hiện của từng mặt.
\end{itemize}
}
\end{tcolorbox}
\end{minipage}
\hfill
\begin{minipage}[t]{0.49\textwidth}\vspace{0pt}
\begin{tcolorbox}[
  enhanced,
  colback=teal!4!white,
  colframe=teal!75!black,
  arc=2mm, boxrule=0.8pt,
  left=2mm, right=2mm, top=1.8mm, bottom=1.8mm,
  title={\footnotesize\bfseries\color{white}\faSliders\ 2. Đọc và nhận xét kết quả}
]
{\footnotesize
\begin{itemize}\setlength{\itemsep}{1.5pt}\setlength{\leftskip}{-2.5mm}
  \item Nhận biết mặt xuất hiện \textbf{nhiều lần nhất} (có số vạch kiểm đếm lớn nhất).
  \item Nhận biết mặt xuất hiện \textbf{ít lần nhất} (có số vạch kiểm đếm nhỏ nhất).
  \item Nhận xét tính ngẫu nhiên: Kết quả của mỗi lần thử là hoàn toàn độc lập với các lần trước đó.
\end{itemize}
}
\end{tcolorbox}
\end{minipage}

\vspace{-2.5mm}
% --- KHỐI 3: SƠ ĐỒ TƯ DUY TIKZ ---
\noindent
\begin{tcolorbox}[
  enhanced,
  colback=orange!4!white,
  colframe=orange!80!black,
  arc=2mm, boxrule=0.8pt,
  left=2mm, right=2mm, top=1.5mm, bottom=1.5mm,
  title={\footnotesize\bfseries\color{white}\faSitemap\ SƠ ĐỒ TƯ DUY: QUY TRÌNH THỰC HÀNH TRẢI NGHIỆM THỐNG KÊ}
]
\centering
\begin{tikzpicture}[
  box/.style={draw, rounded corners=2pt, align=center, font=\scriptsize\bfseries, inner sep=2.5pt, minimum height=0.75cm},
  arr/.style={-{Stealth[length=2mm]}, thick}
]
  \node[box, fill=purple!15, draw=purple!80!black, text width=2.6cm] (e1) at (0, 0) {THỬ NGHIỆM\\Tung xu, gieo xúc xắc\\Lấy bóng ngẫu nhiên};
  \node[box, fill=teal!15, draw=teal!80!black, text width=2.6cm] (e2) at (3.5, 0) {KIỂM ĐẾM\\Gạch sổ tay từng lần\\Gộp cụm 5 vạch};
  \node[box, fill=blue!12, draw=blue!80!black, text width=2.6cm] (e3) at (7.0, 0) {LẬP BẢNG SỐ LIỆU\\Ghi con số tổng kết\\tương ứng mỗi mặt};
  \node[box, fill=orange!15, draw=orange!80!black, text width=2.6cm] (e4) at (10.5, 0) {KẾT LUẬN\\Mặt nhiều nhất\\Mặt ít nhất};

  \draw[arr, purple!80!black] (e1) -- (e2);
  \draw[arr, teal!80!black] (e2) -- (e3);
  \draw[arr, blue!80!black] (e3) -- (e4);
\end{tikzpicture}
\end{tcolorbox}

\vspace{-2.5mm}
% --- KHỐI 4: VÍ DỤ & MẸO TRÁNH BẪY ---
\noindent
\begin{minipage}[t]{0.49\textwidth}\vspace{0pt}
\begin{tcolorbox}[
  enhanced,
  colback=green!4!white,
  colframe=green!60!black,
  arc=2mm, boxrule=0.8pt,
  left=2mm, right=2mm, top=1.8mm, bottom=1.8mm,
  title={\footnotesize\bfseries\color{white}\faCheckCircle\ Ví dụ thực hành chuẩn mực}
]
{\footnotesize
\textbf{Gieo xúc xắc 20 lần, ghi nhận số lần xuất hiện:}\\
\begin{center}
\renewcommand{\arraystretch}{1.1}
\begin{tabular}{|l|c|c|c|c|c|c|}
  \hline
  \textbf{Mặt (chấm)} & $1$ & $2$ & $3$ & $4$ & $5$ & $6$ \\
  \hline
  \textbf{Số lần} & $3$ & $4$ & $5$ & $2$ & $4$ & $2$ \\
  \hline
\end{tabular}
\end{center}
-- Tổng số lần gieo: $3 + 4 + 5 + 2 + 4 + 2 = \mathbf{20\text{ lần}}$.\\
-- Mặt xuất hiện nhiều nhất: \textbf{Mặt 3 chấm} ($5$ lần).\\
-- Mặt xuất hiện ít nhất: \textbf{Mặt 4 và mặt 6 chấm} (cùng $2$ lần).
}
\end{tcolorbox}
\end{minipage}
\hfill
\begin{minipage}[t]{0.49\textwidth}\vspace{0pt}
\begin{tcolorbox}[
  enhanced,
  colback=red!4!white,
  colframe=red!70!black,
  arc=2mm, boxrule=0.8pt,
  left=2mm, right=2mm, top=1.8mm, bottom=1.8mm,
  title={\footnotesize\bfseries\color{white}\faExclamationTriangle\ Mẹo tránh bẫy trắc nghiệm}
]
{\footnotesize
\begin{itemize}\setlength{\itemsep}{1.5pt}\setlength{\leftskip}{-2.5mm}
  \item \textbf{Bẫy đếm thiếu hoặc thừa số lần thử:}\\
  Sau khi hoàn thành bảng số liệu, luôn cộng tổng số lần của các mặt lại để đối chiếu xem có đúng bằng tổng số lần thử ban đầu của đề bài không!
  \item Không được bỏ sót các mặt có số lần xuất hiện bằng $0$.
\end{itemize}
}
\end{tcolorbox}
\end{minipage}


\end{document}
